\begin{definition}
Let $0<\beta\le1$, let $\Omega,\Omega_{\mathrm{cl}}\subseteq\R^n$, let
$u\colon\R^n\to\R$ and let $g\colon\R^n\to\R^n$.  We write
$u\in C^{1+\beta}(\Omega,\Omega_{\mathrm{cl}})$ \emph{with gradient $g$},
denoted
$\mathrm{MemC1HolderOn}(n,\beta,\Omega,\Omega_{\mathrm{cl}},u,g)$, if
\begin{itemize}
\item[(a)] $u$ is differentiable at every point of $\Omega$, with
  $\nabla u=g$ there;
\item[(b)] $g$ is continuous on $\Omega_{\mathrm{cl}}$; and
\item[(c)] there is $C\ge0$ with
  $|g(x)-g(y)|\le C\,|x-y|^{\beta}$ for all
  $x,y\in\Omega_{\mathrm{cl}}$.
\end{itemize}

\noindent\textbf{The two sets, and why the class is written on two of them.}
Taking $\Omega_{\mathrm{cl}}=\Omega$ gives back the one-set reading of the class
verbatim --- the three clauses then range over the same set --- so this
definition is a refinement and not a change of meaning on the diagonal.  The
instance used in \S4 and \S5 is $\Omega_{\mathrm{cl}}=\overline{\Omega}$, and
that instance \emph{is} the paper's class $C^{1+\beta}=C^{1,\beta}$: the Note's
Definitions paragraph at l.~278 says that ``$C^{\eta}$ and
$C^{1+\beta}=C^{1,\beta}$ are the H\"older classes on $\overline{\Omega'}$'',
which for $C^{1,\beta}$ is a statement about the \emph{gradient's} H\"older
seminorm, namely clauses (b) and (c) above; while the two conclusions being
formalised are displayed at l.~1352 and l.~1574 as
$u\in C^{1+\beta}(\Omega')$, on the \emph{open} $\Omega'$, which is clause~(a).
The paper writes exactly this split in its own notation at l.~559, where the
conclusion of Lemma~3.1 --- the cited result (C2) --- is
$u\in C^{1}(B_R)\cap C^{\beta}(\overline B_R)$: differentiability on the open
ball, H\"older control on the closed one, as an intersection of two classes over
two different sets.  Theorem~7's conclusion at l.~1635 is likewise
$F\in C^{1+\beta}_{\mathrm{loc}}(\Omega)$ on an open set.

\noindent\textbf{Why clause (a) must not be asked on
$\Omega_{\mathrm{cl}}=\overline{\Omega}$.}  Differentiability at a point is a
two-sided condition in a full neighbourhood of that point in $\R^n$, so asking
clause (a) at a boundary point of $\Omega$ constrains $u$ outside
$\overline{\Omega}$ as well, and on the data the consumers supply it is not
merely unproven but false.  Let $K\subseteq(-1,1)$ be a nowhere dense closed set
of positive measure, $S=\{x\in\R^n:x_1\in K\}$, $\Omega=B_2(0)$,
$\Omega'=B_1(0)\setminus S$ and $u(x)=\bigl|[0,x_1]\cap K\bigr|$.  Then
$\Omega'$ is open, $u$ is Lipschitz with weak gradient $0$ almost everywhere on
$\Omega'$, so $g\equiv0$ is a continuous $\beta$-H\"older representative on
every set, and clause (a) holds on $\Omega'$ with $u$ itself as the
representative --- at a point with $x_1\notin K$ a whole neighbourhood of $x_1$
misses the closed set $K$, so $u$ is locally constant there; but $\overline{\Omega'}=\overline B_1(0)$, and a function
differentiable at every point of $\overline B_1(0)$ with gradient $0$ there is
constant on the convex set $B_1(0)$, whereas $u$ increases by $|K|>0$ across it.
So no function agreeing with $u$ almost everywhere can satisfy clause (a) on
$\overline{\Omega'}$.  The same example rules out
$\mathrm{interior}\,\overline{\Omega'}$, which here is $B_1(0)$.
\end{definition}
